%% ============================================================
%%  サンプル章
%%  このテンプレートが扱える書き方を一通り含んでいる。
%%  自分のノートを書き始めるときは、中身を差し替えて使う。
%% ============================================================

\chapter{距離空間と収束}
\label{ch:introduction}

本章では距離空間を定義し、点列の収束と完備性を扱う。
集合と写像に関する前提は付録にまとめてある。

\section{距離空間}
\label{sec:metric-space}

\begin{definition}{距離空間}{metric-space}
集合 $M$ と写像 $d \colon M \times M \to \R$ の組 $(M, d)$ が
\textbf{距離空間}であるとは、任意の $x, y, z \in M$ に対して次が成り立つことをいう。
\begin{enumerate}
  \item $d(x, y) \geq 0$ であり、$d(x, y) = 0$ と $x = y$ は同値である（正値性）。
  \item $d(x, y) = d(y, x)$（対称性）。
  \item $d(x, z) \leq d(x, y) + d(y, z)$（三角不等式）。
\end{enumerate}
このとき $d$ を $M$ 上の\textbf{距離}という。
\end{definition}

\begin{example}{Euclid 距離}{euclidean}
$M = \R^n$ とし、
\[
  d(x, y) = \norm{x - y} = \left( \sum_{i=1}^{n} \abs{x_i - y_i}^2 \right)^{1/2}
\]
と定めると、$(\R^n, d)$ は距離空間である。三角不等式は Cauchy--Schwarz の不等式から従う。
\end{example}

\begin{remark}{距離の値域}{codomain}
定義では $d$ の値域を $\R$ としたが、正値性から実際には $d(M \times M) \subset [0, \infty)$ である。
$\infty$ を値に許す一般化（拡張距離）はここでは扱わない。
\end{remark}

\section{収束と極限}
\label{sec:convergence}

\begin{definition}{点列の収束}{convergence}
距離空間 $(M, d)$ の点列 $(x_n)_{n \in \N}$ が $x \in M$ に\textbf{収束する}とは、
\[
  \forall \varepsilon > 0,\ \exists N \in \N,\ \forall n \geq N,\ d(x_n, x) < \varepsilon
\]
が成り立つことをいう。このとき $x$ を $(x_n)$ の\textbf{極限}といい、
$x_n \to x$ と書く。
\end{definition}

\begin{proposition}{極限の一意性}{limit-unique}
距離空間 $(M, d)$ の点列 $(x_n)$ が $x$ にも $y$ にも収束するならば、$x = y$ である。
\end{proposition}

\begin{proof}
$x \neq y$ と仮定する。距離の正値性（定義~\ref{def:metric-space}）より $d(x, y) > 0$ であるから、
$\varepsilon = d(x, y) / 2 > 0$ とおける。収束の定義~\ref{def:convergence}より、
ある $N_1, N_2 \in \N$ が存在して
\[
  n \geq N_1 \implies d(x_n, x) < \varepsilon,
  \qquad
  n \geq N_2 \implies d(x_n, y) < \varepsilon
\]
が成り立つ。$n = \max\{N_1, N_2\}$ をとると、三角不等式より
\begin{align*}
  d(x, y)
  &\leq d(x, x_n) + d(x_n, y) \\
  &< \varepsilon + \varepsilon \\
  &= d(x, y)
\end{align*}
となり矛盾する。よって $x = y$ である。
\end{proof}

\begin{lemma}{収束列は有界}{convergent-bounded}
距離空間 $(M, d)$ の収束列 $(x_n)$ に対し、ある $x \in M$ と $R > 0$ が存在して、
すべての $n \in \N$ について $d(x_n, x) \leq R$ が成り立つ。
\end{lemma}

\begin{proof}
$x_n \to x$ とする。$\varepsilon = 1$ に対して $N \in \N$ をとれば、
$n \geq N$ のとき $d(x_n, x) < 1$ である。残りは有限個なので
\[
  R = \max\{1,\ d(x_1, x),\ \ldots,\ d(x_{N-1}, x)\}
\]
とおけばよい。
\end{proof}

\section{完備性}
\label{sec:completeness}

\begin{definition}{Cauchy 列}{cauchy}
距離空間 $(M, d)$ の点列 $(x_n)$ が \textbf{Cauchy 列}であるとは、
\[
  \forall \varepsilon > 0,\ \exists N \in \N,\ \forall m, n \geq N,\ d(x_m, x_n) < \varepsilon
\]
が成り立つことをいう。
\end{definition}

\begin{claim}{収束列は Cauchy 列}{convergent-cauchy}
距離空間の収束列は Cauchy 列である。
\end{claim}

\begin{proof}
$x_n \to x$ とし、$\varepsilon > 0$ をとる。定義~\ref{def:convergence}より、
$n \geq N$ ならば $d(x_n, x) < \varepsilon / 2$ となる $N$ が存在する。
$m, n \geq N$ のとき、三角不等式より
\[
  d(x_m, x_n) \leq d(x_m, x) + d(x, x_n) < \frac{\varepsilon}{2} + \frac{\varepsilon}{2} = \varepsilon
\]
である。
\end{proof}

\begin{definition}{完備距離空間}{complete}
距離空間 $(M, d)$ が\textbf{完備}であるとは、$M$ の任意の Cauchy 列が $M$ の点に収束することをいう。
\end{definition}

逆向きの主張、すなわち Cauchy 列が必ず収束することは一般には成り立たない。

\begin{example}{完備でない距離空間}{incomplete}
$M = \Q$ に Euclid 距離を入れた空間は完備ではない。
$\sqrt{2}$ に収束する有理数列は $\Q$ の中では極限を持たないが、Cauchy 列である。
\end{example}

\begin{theorem}{閉部分集合の完備性}{closed-complete}
完備距離空間 $(M, d)$ の部分集合 $A \subset M$ が閉集合であれば、
$(A, d)$ も完備距離空間である。
\end{theorem}

\begin{proof}
$(x_n)$ を $A$ の Cauchy 列とする。これは $M$ の Cauchy 列でもあるから、
$M$ の完備性よりある $x \in M$ に収束する。$A$ は閉集合であり
$x$ は $A$ の点列の極限であるから $x \in A$ である。
よって $(x_n)$ は $A$ の点に収束する。
\end{proof}

\begin{corollary}{有界閉区間の完備性}{closed-interval}
$a \leq b$ なる実数に対し、$([a, b], \abs{\cdot})$ は完備距離空間である。
\end{corollary}

\begin{proof}
$\R$ は完備であり、$[a, b]$ は $\R$ の閉集合である。定理~\ref{thm:closed-complete}を適用すればよい。
\end{proof}

\begin{memo*}
完備性は距離に依存する性質であって、位相だけでは決まらない。
開区間 $(0, 1)$ は $\R$ と同相だが、Euclid 距離のもとでは完備でない。
\end{memo*}

\section{アルゴリズムと図}
\label{sec:extras}

反復で不動点を求める手続きは、次のように書ける。

\begin{algorithm}{alg:fixed-point}
入力: 縮小写像 $T \colon M \to M$、初期点 $x_0 \in M$、許容誤差 $\varepsilon > 0$ \\
$n \leftarrow 0$ \\
$d(x_{n+1}, x_n) < \varepsilon$ となるまで $x_{n+1} \leftarrow T(x_n)$ と更新し、$n$ を 1 増やす \\
出力: $x_n$
\end{algorithm}

図は tikz で書く。PDF には出るが、サイトでは省略される。

\begin{figure}[H]
  \centering
  \begin{tikzpicture}[scale=1.2]
    \draw[->] (-0.3,0) -- (3.5,0) node[right] {$n$};
    \draw[->] (0,-0.3) -- (0,2.2) node[above] {$d(x_n, x)$};
    \draw[dashed] (0,0.2) -- (3.3,0.2) node[right] {$\varepsilon$};
    \foreach \x/\y in {0.5/1.8, 1.0/1.1, 1.5/0.7, 2.0/0.42, 2.5/0.26, 3.0/0.16}
      \fill (\x,\y) circle (1.6pt);
  \end{tikzpicture}
  \caption{収束の様子。有限回で $\varepsilon$ を下回る。}
  \label{fig:convergence}
\end{figure}

サイトにだけ図を出したいときは、site.config.mjs の demos に登録した名前を
行頭で呼び出す。次の位置に差し込まれる（PDF には何も出ない）。

\demohint{convergence}
